\section{Setting}

\subsection{Two registers of one population}

Artificial objects in Earth orbit are catalogued more than once. We used two
sources, both openly retrievable without credentials, both retrieved on
18~August~2026. CelesTrak's satellite catalogue \citep{celestrak2026satcat}
carried 70,292 objects. The General Catalog of Artificial Space Objects (GCAT)
\citep{mcdowell2020gcat}, maintained independently and published under CC~BY
4.0, is split across many files. The pipeline read four of them: a main
catalogue of 69,999 records covering 69,391 numbered objects, an auxiliary
catalogue of 11,962 records, a catalogue of 1,802 objects that failed to orbit,
and a catalogue of 354 numbered objects above 100,000. Two registers describing
one population is a useful setting because neither has to be treated as ground
truth: where they disagree, the disagreement itself is the finding.

\subsection{The pipeline and its gate}

The pipeline reconciles objects by catalogue number, emits an RDF graph in which
every defect is a typed node carrying its evidence, and validates that graph
against three layers of SHACL shapes \citep{shacl2017}, one shape per defect
class. On the published run it gated seven defect classes, among them entries
one register says correspond to no real object while the other still carries
them, and objects on which the two registers disagree about whether the object
is still in orbit.

Verification was deliberately over-provisioned. Three independent checks ran on
every build: a SHACL processor over all three shape layers, a second RDF engine
written in a different language that loaded and validated the same ontology, and
a dual-computation gate that recomputed every published count twice. The third
is the subject of this paper, and its structure is the obvious one:

\begin{quote}\small
\texttt{python\_side()} reads the source files and returns seven counts.
\texttt{sparql\_side()} parses the emitted graph and issues one
\texttt{SELECT (COUNT(DISTINCT ?n))} per defect class. \texttt{main()} compares
them key by key and calls \texttt{sys.exit(1)} if any pair differs.
\end{quote}

The diversity is real. One path never constructs a graph and the other never
sees the source files. One counts Python set members, the other counts distinct
SPARQL bindings over a materialised graph of 2,330,660 triples, two orders of
magnitude larger than the input. Parsing that graph takes over an hour and
roughly 5~GB of memory, so the gate is not run casually, which is part of why its
verdict carried weight. On the published run it agreed on all seven counts.
